\documentclass[11pt,reqno]{amsart}

\usepackage{amsmath,amssymb,amsthm}
\usepackage{mathtools}
\usepackage[margin=1.1in]{geometry}
\usepackage{hyperref}
\usepackage{enumitem}
\hypersetup{colorlinks=true,linkcolor=blue,citecolor=blue,urlcolor=blue}
\setlength{\emergencystretch}{3em}

\theoremstyle{plain}
\newtheorem{theorem}{Theorem}[section]
\newtheorem{lemma}[theorem]{Lemma}
\newtheorem{proposition}[theorem]{Proposition}
\newtheorem{corollary}[theorem]{Corollary}
\newtheorem{conjecture}[theorem]{Conjecture}
\theoremstyle{definition}
\newtheorem{definition}[theorem]{Definition}
\newtheorem{remark}[theorem]{Remark}

\newcommand{\C}{\mathbb{C}}
\newcommand{\Z}{\mathbb{Z}}
\newcommand{\Q}{\mathbb{Q}}
\newcommand{\N}{\mathbb{N}}
\newcommand{\E}{\mathrm{E}}
\newcommand{\CT}{\operatorname{CT}}
\newcommand{\ch}{\operatorname{ch}}
\newcommand{\wt}{\operatorname{wt}}
\newcommand{\supp}{\operatorname{supp}}
\newcommand{\Gal}{\operatorname{Gal}}
\newcommand{\ord}{\operatorname{ord}}
\newcommand{\degW}{\deg_W}
\DeclareMathOperator{\lc}{lc}

\begin{document}

\title[The Gaussian--Moments and Nullcone conjectures in dimension two]
{A proof of the two--dimensional Gaussian--Moments and Nullcone conjectures,
formally verified}

\author{Eliott Cassidy}
\email{eliottcassidy2000@gmail.com}

\date{\today}

\begin{abstract}
Let $\E$ be the Wick (Gaussian) expectation on $\C[Z,W]$ determined by
$\E[Z^aW^b]=a!\,[a=b]$, the moment functional of a single complex Gaussian.
The \emph{Gaussian--Moments Conjecture} and the \emph{Nullcone Conjecture} of
Derksen, van den Essen and Zhao assert, respectively, that the kernel of the
Gaussian moment functional is a Mathieu--Zhao space, and that a moment--null polynomial has
one--sided charge. The former implies the Jacobian Conjecture but fails in higher dimension, so
dimension two is the sharp surviving case. We prove both conjectures in dimension two: if
$\E(P^m)=0$ for all $m\ge 1$ then $\E(Q\,P^m)=0$ for all large $m$, and every $P$ with $\E(P^m)=0$
for all $m$ has strictly one--sided charge. The proof deduces the moment statement from the
nullcone statement and reduces the residual analytic content to the one--variable
Duistermaat--van der Kallen constant--term theorem. We give a self--contained, valuation--free proof of the latter through a
Galois orbit--product on the splitting field of $X^{M}-tR$, fed by the Weierstrass small--root
product $\Pi=c\,t$, whose analytic heart $\partial_t h(0,t)=0$ is proved inside a single power--series
frame. The entire argument has been formalized in the Lean~4 proof assistant on top of
\texttt{mathlib}; the top--level theorem is machine--checked and depends only on the three
foundational axioms of that library.
\end{abstract}

\maketitle

\section{Introduction}

\subsection{The conjecture}
Fix the polynomial ring $\C[Z,W]$ in two variables and regard $W=\bar Z$ as the conjugate of a single
complex Gaussian variable $Z$. The associated \emph{Wick}, or \emph{Gaussian}, expectation is the
linear functional
\begin{equation}\label{eq:E-def}
  \E\colon \C[Z,W]\longrightarrow \C,\qquad
  \E\big[Z^aW^b\big]=a!\,[a=b],
\end{equation}
where $[a=b]$ is $1$ if $a=b$ and $0$ otherwise. Equivalently, writing $Z=x+iy$, $W=x-iy$,
$\E$ is integration against the standard planar Gaussian, and \eqref{eq:E-def} is the classical Wick
formula for its moments. The functional $\E$ annihilates every monomial off the diagonal, i.e. every
monomial of nonzero \emph{charge}
\[
  \ch(Z^aW^b)\;=\;a-b .
\]

A linear subspace $\mathcal M$ of a commutative $\C$--algebra $A$ is a \emph{Mathieu--Zhao space}
(Zhao \cite{Zhao2010}; originally \emph{Mathieu subspace}) if, whenever $a\in A$ satisfies $a^m\in\mathcal M$
for all $m\ge1$, one has $b\,a^m\in\mathcal M$ for all sufficiently large $m$ and every $b\in A$. These
spaces isolate the algebraic content of a family of conjectures descending from the Jacobian
Conjecture \cite{vandenEssen2000}: Mathieu's conjecture on integrals of $G$--finite functions
\cite{Mathieu1997}, Zhao's Image Conjecture and its Hessian--nilpotent form
\cite{ZhaoImage,Zhao2007a}, and their common Gaussian shadow.

Derksen, van den Essen and Zhao \cite{DerksenEssenZhao2017} formulated, and connected to the Jacobian
Conjecture, two conjectures about the Gaussian moment functional. The first asks that its kernel be a
Mathieu--Zhao space.

\begin{conjecture}[Gaussian--Moments Conjecture, {\cite{DerksenEssenZhao2017}}]\label{conj:GMC}
The kernel of the Gaussian moment expectation is a Mathieu--Zhao space of the ambient polynomial ring.
\end{conjecture}

They proved that Conjecture~\ref{conj:GMC} implies the Jacobian Conjecture, and observed that it does
not extend to arbitrary dimension: a counterexample to the unrestricted moment--vanishing statement
shows that the Gaussian--Moments phenomenon fails once enough Gaussian variables are present. In the
two--variable ring $\C[Z,W]$, which realizes a single complex Gaussian $Z$ with $W=\bar Z$ via
\eqref{eq:E-def}, the conjecture survives; this is the sharp case, and it is our first theorem.

\begin{theorem}[Gaussian--Moments Conjecture in dimension two]\label{thm:gmc2}
Let $P,Q\in\C[Z,W]$. If $\E(P^m)=0$ for all $m\ge 1$, then there is $N$ such that $\E(Q\,P^m)=0$
for all $m\ge N$.
\end{theorem}

The companion \emph{Nullcone Conjecture} of \cite{DerksenEssenZhao2017} makes the underlying instability
structural: a polynomial all of whose central Gaussian moments vanish should lie in the nullcone of the
relevant torus action, i.e. be destabilized by a one--parameter subgroup. In dimension two the torus is
the single hyperbolic $\C^\times$ acting with weight $\ch(s)=s_0-s_1$ on $Z^{s_0}W^{s_1}$, and its
nullcone consists precisely of the strictly one--sided polynomials. Our second theorem is this
two--dimensional Nullcone Conjecture.

\begin{theorem}[Nullcone Conjecture in dimension two]\label{thm:nc2}
Let $P\in\C[Z,W]$ and write $\supp P\subset\N^2$ for its exponent support. If $\E(P^m)=0$ for all
$m\ge 1$, then $P$ has one--sided charge:
either $\ch(s)\ge 1$ for every $s\in\supp P$, or $\ch(s)\le -1$ for every $s\in\supp P$.
\end{theorem}

We abbreviate Theorems~\ref{thm:gmc2} and \ref{thm:nc2} by $\mathrm{GMC}(2)$ and $\mathrm{NC}(2)$.
Theorem~\ref{thm:nc2} is the structural heart of Theorem~\ref{thm:gmc2}: it says that the only way for
all central powers of $P$ to be moment--null is for $P$ to be strictly one--sided (no monomial can have
charge $0$, and all charges must agree in sign). The passage $\mathrm{NC}(2)\Rightarrow\mathrm{GMC}(2)$
is elementary (Section~\ref{sec:charge}); the content is Theorem~\ref{thm:nc2} itself.

\subsection{Method}
The proof of Theorem~\ref{thm:nc2} confronts one, and only one, genuinely analytic input: the
one--variable \emph{Duistermaat--van der Kallen theorem} on constant terms in powers of a Laurent
polynomial \cite{DvdK1998}. In the exact form we need it reads:

\begin{theorem}[Duistermaat--van der Kallen, finite--support form]\label{thm:dvdk}
Let $f(z)=\sum_{i} c_i z^{q_i}$ be a Laurent polynomial with distinct exponents $q_i\in\Z$ and nonzero
coefficients $c_i\in\C$, whose support \emph{straddles zero}: either $0$ occurs among the $q_i$, or both a
negative and a positive exponent occur. Then there is $m\ge 1$ with $\CT(f^m)\ne 0$, where $\CT$ denotes
the constant Laurent coefficient.
\end{theorem}

The classical theorem of Duistermaat and van der Kallen \cite{DvdK1998} classifies the Laurent
polynomials all of whose positive powers have vanishing constant term; Theorem~\ref{thm:dvdk} is the
contrapositive in the two--sided case, together with the trivial zero--charge case. We do not import
the analytic proof of \cite{DvdK1998}. Instead we give a self--contained algebraic proof
(Sections~\ref{sec:univariate}--\ref{sec:analytic}) that runs as follows. Clearing denominators turns
$f$ into a genuine polynomial and reduces the vanishing of all $\CT(f^m)$ to the vanishing of the
central coefficients $D_m=[X^{Mm}]R(X)^m$ of a single polynomial $R$; these are precisely the
coefficients governing the small roots of
\[
  \Phi \;=\; X^{M}-t\,R(X)\ \in\ \C(t)[X].
\]
Over the splitting field of $\Phi$, a transitive Galois action and a Vieta identity force a
contradiction \emph{as soon as} the product of the $M$ small roots equals $c\,t$ with $c\ne0$
(Section~\ref{sec:orbit}). That product identity is exactly the Weierstrass small--root product
$\Pi=c\,t$, whose only nontrivial ingredient is the scalar analytic identity
$\partial_t h(0,t)=0$ under $D_m\equiv 0$, proved inside the single power--series frame
$(\C(\!(x)\!))[\![t]\!]$ (Section~\ref{sec:analytic}). No extension of a valuation to an algebraic
closure, and no Puiseux expansion, is required.

\subsection{Formalization}
Every step above --- Theorems~\ref{thm:gmc2}, \ref{thm:nc2}, \ref{thm:dvdk}, the univariate reduction,
the orbit--product contradiction, and the Weierstrass analysis --- has been formalized in the Lean~4
proof assistant on top of \texttt{mathlib} \cite{mathlib}. The top--level statement is
\begin{quote}\ttfamily
theorem GMC2.gmc2 (P Q : MvPolynomial (Fin 2) $\C$)\\
\hspace*{2em}(hnull : $\forall$ m, 1 $\le$ m $\to$ E (P \^{} m) = 0) :\\
\hspace*{2em}$\exists$ N, $\forall$ m $\ge$ N, E (Q * P \^{} m) = 0
\end{quote}
and a kernel audit reports
\[
  \texttt{\#print axioms GMC2.gmc2} \;=\; [\,\mathtt{propext},\ \mathtt{Classical.choice},\ \mathtt{Quot.sound}\,],
\]
the three standard axioms of \texttt{mathlib}: there is no \texttt{sorry}, no appeal to
\texttt{native\_decide}, and no additional axiom anywhere in the transitive dependency. In particular
the one--variable input, Theorem~\ref{thm:dvdk}, is itself formalized rather than assumed. The proof
is described module by module in Section~\ref{sec:lean}.

\subsection{Organization}
Section~\ref{sec:charge} records the functional $\E$ and the elementary reduction
$\text{NC}(2)\Rightarrow\text{GMC}(2)$. Section~\ref{sec:nc2} reduces $\text{NC}(2)$ to
Theorem~\ref{thm:dvdk}. Section~\ref{sec:univariate} reduces Theorem~\ref{thm:dvdk} to a single
polynomial statement (the \emph{small--root packet product}). Section~\ref{sec:orbit} closes that
statement by a Galois orbit--product. Section~\ref{sec:analytic} supplies the Weierstrass small--root
product and its analytic core. Section~\ref{sec:lean} documents the Lean formalization and its
verification. Section~\ref{sec:remarks} collects remarks on sharpness and on the Jacobian Conjecture.

\section{The Wick functional and the charge reduction}\label{sec:charge}

Write monomials of $\C[Z,W]$ as $Z^{s_0}W^{s_1}$ with $s=(s_0,s_1)\in\N^2$, and set
\[
  \wt(s)=\begin{cases} s_0! & s_0=s_1,\\ 0 & \text{otherwise,}\end{cases}
  \qquad
  \E(P)=\sum_{s\in\supp P} P_s\,\wt(s),
\]
where $P_s$ is the coefficient of $Z^{s_0}W^{s_1}$ in $P$. Thus $\E$ agrees with \eqref{eq:E-def} and
annihilates every monomial of nonzero charge $\ch(s)=s_0-s_1$.

\begin{lemma}\label{lem:charge-kill}
If $P,R\in\C[Z,W]$ have the property that every monomial of $P$ and every monomial of $R$ has charge
$\ge 1$ (respectively $\le -1$), then $\E(R)=0$; more generally $\E(PR)=0$ whenever every monomial of
$PR$ has nonzero charge.
\end{lemma}

\begin{proof}
Charge is additive under multiplication, $\ch(s+t)=\ch(s)+\ch(t)$, so every monomial of a product of
strictly one--sided polynomials has charge of the same strict sign, hence nonzero; and $\wt$ vanishes off
charge $0$.
\end{proof}

\begin{proposition}[$\text{NC}(2)\Rightarrow\text{GMC}(2)$]\label{prop:nc2-to-gmc2}
Assume Theorem~\ref{thm:nc2}. Then Theorem~\ref{thm:gmc2} holds; explicitly, if $\E(P^m)=0$ for all
$m\ge1$ then $\E(Q\,P^m)=0$ for all $m\ge \deg_W Q+1$ (in the case of positive charge; symmetrically for
negative charge).
\end{proposition}

\begin{proof}
By Theorem~\ref{thm:nc2}, $P$ is strictly one--sided; say every monomial of $P$ has charge $\ge1$
(the other case is symmetric under $Z\leftrightarrow W$). Every monomial of $P^m$ then has charge
$\ge m$. A monomial $Z^{a}W^{b}$ of $Q$ has $b\le \deg_W Q$, so its charge is $\ge -\deg_W Q$; hence a
monomial of $Q\,P^m$ has charge $\ge m-\deg_W Q\ge 1$ once $m\ge \deg_W Q+1$. By
Lemma~\ref{lem:charge-kill}, $\E(Q\,P^m)=0$ for all such $m$.
\end{proof}

The threshold $N=\deg_W Q+1$ is explicit and uniform: the whole difficulty is concentrated in
Theorem~\ref{thm:nc2}, to which we now turn.

\section{Reduction of NC$(2)$ to the constant--term theorem}\label{sec:nc2}

Suppose $\E(P^m)=0$ for all $m\ge1$ but, for contradiction, that $P$ is \emph{not} one--sided: its
charge support is not contained in $\Z_{\ge1}$ nor in $\Z_{\le-1}$. Because $\E$ only sees the diagonal,
only the \emph{charge--balanced} part of each power matters:
\begin{equation}\label{eq:wick}
  \E(P^m)\;=\;\sum_{\substack{r\in\N^{\,\supp P}\\ |r|=m,\ \sum_i r_i\,\ch(s_i)=0}}
      \binom{m}{r}\Big(\prod_i P_{s_i}^{\,r_i}\Big)\,\big(\text{diagonal Wick weight}\big),
\end{equation}
the sum over compositions $r$ of $m$ across the monomials $s_i$ of $P$ whose total charge is zero.
Passing to the associated one--dimensional charge data $q_i=\ch(s_i)\in\Z$ with coefficients
$c_i=P_{s_i}\in\C^\times$ isolates the \emph{angular} content of \eqref{eq:wick}: the balanced
compositions of $q=(q_i)$ are exactly the constant--term channels of the Laurent polynomial
$f(z)=\sum_i c_i z^{q_i}$. The radial content --- the factorials $\wt$ and the Frobenius amplification
that upgrades a single balanced face to all of them --- is carried by an elementary but non--vacuous
combinatorial argument (the lowest--balanced--face / Kummer--carry mechanism, formalized in the module
\texttt{GMC2NC2}) whose only external call is the following: the constant--term relations of $f$ cannot
all vanish.

Concretely, let $q\colon\{1,\dots,k\}\to\Z$ be injective with values the distinct charges and
$c_i\ne0$ the corresponding coefficients. For $m\ge0$ define the universal constant--term polynomial
\begin{equation}\label{eq:ctr}
  \mathcal C_m(q)\;=\;
  \sum_{\substack{r\in\N^k,\ |r|=m\\ \sum_i r_i q_i=0}} \binom{m}{r}\, X^{r}\ \in\ \Q[X_1,\dots,X_k],
\end{equation}
so that $\mathcal C_m(q)(c)=\CT(f^m)$. The hypothesis $\E(P^m)=0$ for all $m$ forces, after the
Frobenius reduction, $\mathcal C_m(q)(c)=0$ for all $m\ge1$. Since $P$ is not one--sided, the charge set
$\{q_i\}$ straddles zero. Theorem~\ref{thm:dvdk} then provides $m\ge1$ with $\mathcal C_m(q)(c)\ne0$, a
contradiction. Hence $P$ is one--sided, proving Theorem~\ref{thm:nc2}. \qed

\begin{remark}
The zero--charge case of Theorem~\ref{thm:dvdk} is immediate: if some $q_{i_0}=0$, then at $m=1$ the
single composition supported at $i_0$ is the unique balanced channel and $\mathcal C_1(q)(c)=c_{i_0}\ne0$.
The content is therefore the two--sided, no--zero case, treated next.
\end{remark}

\section{The univariate reduction: from constant terms to a small--root packet}\label{sec:univariate}

Fix injective $q$ straddling zero with no zero charge, and coefficients $c_i\ne0$. Let
$M=-\min_i q_i>0$ and form the genuine polynomial
\[
  R(X)\;=\;\sum_i c_i\,X^{\,q_i+M}\ \in\ \C[X].
\]
The shift is chosen so that $R(0)=c_{i_0}\ne0$ (the unique minimal charge) and
$M<\deg R$ (the unique maximal charge is strictly positive). A direct expansion identifies the
constant--term relations \eqref{eq:ctr} with central coefficients of powers of $R$:
\begin{equation}\label{eq:checkA}
  \big[X^{Mm}\big]\,R(X)^m \;=\; \mathcal C_m(q)(c)\qquad (m\ge0).
\end{equation}
Thus assuming $\mathcal C_m(q)(c)=0$ for all $m\ge1$ is the same as assuming
$D_m:=[X^{Mm}]R^m=0$ for all $m\ge1$.

These central coefficients are the moments of the small roots of
\begin{equation}\label{eq:Phi}
  \Phi\;=\;X^{M}-t\,R(X)\ \in\ \C(t)[X],
\end{equation}
a polynomial of degree $N=\deg R>M$ in $X$ over the rational function field $\C(t)$. Modulo $t$ it
degenerates to $X^{M}$, so exactly $M$ of its $N$ roots are ``small'' (they tend to $0$ as $t\to0$);
their product carries a $t$--adic valuation equal to $1$ when the moments $D_m$ vanish. The whole of
Theorem~\ref{thm:dvdk} reduces to the following single statement about $\Phi$.

\begin{proposition}[Small--root packet product]\label{prop:crux}
Let $R\in\C[X]$ with $1\le M<\deg R$ and $R(0)\ne0$, and suppose $D_m=[X^{Mm}]R^m=0$ for all $m\ge1$.
Then, in the splitting field $L$ of $\Phi=X^M-tR$ over $\C(t)$, there is a finite set $S$ of roots of
$\Phi$ whose product equals $c\,t$ for some $c\in\C^\times$:
\[
  \prod_{\beta\in S}\beta \;=\; c\,t\ \in\ \C(t)\subset L.
\]
\end{proposition}

Granting Proposition~\ref{prop:crux}, Theorem~\ref{thm:dvdk} follows: were $D_m=0$ for all $m\ge1$, the
packet product would be $c\,t$, and Section~\ref{sec:orbit} shows this is impossible; hence some
$D_m\ne0$, i.e. some $\mathcal C_m(q)(c)\ne0$.

\section{The Galois orbit--product contradiction}\label{sec:orbit}

We first record why a packet product equal to $c\,t$ is impossible, independently of \emph{how} the
packet arises. This is the algebraic core, an orbit--product identity for the Galois action on the
roots.

\begin{lemma}\label{lem:irred}
$\Phi=X^M-tR$ is irreducible over $\C(t)$ whenever $1\le M$ and $R(0)\ne0$.
\end{lemma}

\begin{proof}[Sketch]
$\Phi$ is linear in $t$ with content $1$; the coefficients $X^M$ and $-R$ of $t^0,t^1$ are coprime in
$\C[X]$ because $R(0)\ne0$ (so $X\nmid R$). Hence $\Phi$ is irreducible in $\C[X][t]=\C[t][X]$, and by
Gauss's lemma over the fraction field $\C(t)$. (In the formalization this is carried out through the
Bivariate swap $\C[t][X]\cong\C[X][t]$ and \texttt{mathlib}'s primitive/Gauss machinery.)
\end{proof}

Let $L$ be the splitting field of $\Phi$ over $F:=\C(t)$ and $G=\Gal(L/F)$. Since $\Phi$ is irreducible,
$G$ acts transitively on the set $\Omega$ of its roots. The key point is that the orbit--product
identity below holds for an \emph{arbitrary} $G$--stable factor of the product --- one need not know
that $S$ is the ``valuation--positive'' packet, only that its product is $G$--fixed.

\begin{lemma}[Orbit--product]\label{lem:orbit}
Let $G$ act transitively on a finite set $\Omega$ with an equivariant embedding $f\colon\Omega\to L$,
and let $S\subseteq\Omega$ be any subset such that $\pi:=\prod_{\beta\in S}f(\beta)$ is $G$--fixed. Then
for a basepoint $x\in\Omega$,
\[
  \pi^{\,|G|}\;=\;\Big(\prod_{\alpha\in\Omega} f(\alpha)\Big)^{\,|S|\cdot|\mathrm{Stab}_G(x)|}.
\]
\end{lemma}

\begin{proof}
Applying every $g\in G$ and multiplying, $\prod_{g\in G} g\!\cdot\!\pi=\pi^{|G|}$ since $\pi$ is fixed.
On the other hand $\prod_{g\in G} g\!\cdot\!\pi=\prod_{\beta\in S}\prod_{g\in G} f(g\beta)$, and by
transitivity $\prod_{g\in G} f(g\beta)=\big(\prod_{\alpha\in\Omega}f(\alpha)\big)^{|\mathrm{Stab}_G(x)|}$
for every $\beta$. Taking the product over the $|S|$ elements of $S$ gives the claim.
\end{proof}

\begin{proposition}[No packet has product $c\,t$]\label{prop:no-ct}
With $F=\C(t)$, $L$, $G$ as above, suppose $S\subseteq\Omega$ has $G$--fixed product
$\prod_{\beta\in S}\beta=c\,t$ with $c\in\C^\times$. Then a contradiction ensues.
\end{proposition}

\begin{proof}
Because $c\,t\in F$ lies in the base field, it is automatically $G$--fixed, so Lemma~\ref{lem:orbit}
applies with $\pi=c\,t$. The full product of the roots is, by Vieta on \eqref{eq:Phi},
\[
  \prod_{\alpha\in\Omega}\alpha \;=\; (-1)^{N}\,\frac{[X^0]\Phi}{\lc_X\Phi}
   \;=\;(-1)^{N}\,\frac{-t\,R(0)}{-t\,\lc R}
   \;=\;(-1)^{N}\,\frac{R(0)}{\lc R}\;=:\;d\ \in\ \C^\times,
\]
a nonzero \emph{constant}: the factor $t$ cancels between the constant and leading coefficients of
$\Phi$. Substituting into Lemma~\ref{lem:orbit},
\[
  (c\,t)^{\,|G|}\;=\;d^{\,|S|\cdot|\mathrm{Stab}_G(x)|}\quad\text{in }F=\C(t).
\]
The left side has $t$--adic valuation $|G|\ge1$; the right side is a nonzero constant of valuation $0$.
A degree--$|G|$ monomial in $t$ cannot equal a constant, so this is absurd.
\end{proof}

Proposition~\ref{prop:no-ct} is what makes the reduction \emph{valuation--free}. Nothing about the size
of the roots is used; the only inputs are transitivity (Lemma~\ref{lem:irred}), Vieta, and the
elementary fact that a monomial is not a constant. It remains to produce a packet with product $c\,t$,
i.e. to prove Proposition~\ref{prop:crux}. This is where the Weierstrass analysis enters, and it is the
only place any analysis is used.

\section{The Weierstrass small--root product and its analytic core}\label{sec:analytic}

Regard $\Phi$ as a polynomial in $X$ over the complete local ring $A=\C[\![t]\!]$:
\[
  \Phi \;=\; X^{M}-t\,R(X)\ \in\ A[X]\subset A[\![X]\!].
\]
Its reduction mod $t$ is $X^{M}$, so Weierstrass preparation over $A[\![X]\!]$ yields a factorization
\begin{equation}\label{eq:weier}
  \Phi \;=\; P\cdot h,\qquad P\in A[X]\ \text{monic of degree }M,\ P\equiv X^{M}\ (\mathrm{mod}\ t),\quad
  h\in A[\![X]\!]^{\times}.
\end{equation}
$P$ is the \emph{distinguished} (small--root) factor; its $M$ roots are the small roots of $\Phi$, and
the small--root product is $\Pi=(-1)^{M}P.\!\operatorname{coeff}_0$.

\subsection{From power series to a polynomial factorization}
The factorization \eqref{eq:weier} lives in the power--series ring; to feed the packet into
Section~\ref{sec:orbit} we need $P$ to divide $\Phi$ as \emph{polynomials}, so that the roots of $P$ are
genuinely roots of $\Phi$ in $L$.

\begin{lemma}\label{lem:dvd}
$P$ divides $\Phi$ in the polynomial ring $A[X]$.
\end{lemma}

\begin{proof}
Since $P$ is monic, ordinary polynomial division gives $\Phi=P\cdot Q+r$ with $\deg_X r<M$. Coercing to
$A[\![X]\!]$ and comparing with \eqref{eq:weier}, both $(h,0)$ and $(Q,r)$ are Weierstrass divisions of
$\Phi$ by $P$; since $P\equiv X^M\pmod t$ makes $P$ a valid Weierstrass divisor over the Hausdorff ring
$A$, uniqueness of Weierstrass division forces the remainder $r$ to vanish. Hence $P\mid\Phi$ in $A[X]$.
\end{proof}

Mapping $A=\C[\![t]\!]\hookrightarrow\C(\!(t)\!)$ and then into $\Omega=\overline{\C(\!(t)\!)}$, and using
that the polynomial \eqref{eq:Phi} over $\C(t)$ coincides with its image over $\C(\!(t)\!)$, one obtains
$P\mid\Phi$ over $\Omega$. The splitting field $L$ of $\Phi$ over $\C(t)$ embeds $\C(t)$--linearly into
$\Omega$ by an embedding $\psi\colon L\hookrightarrow\Omega$ (both fields split $\Phi$), and the packet
\[
  S=\{\,\beta\in L: \Phi(\beta)=0\ \text{and}\ \psi(\beta)\ \text{is a root of }P\,\}
\]
has, by Vieta on the monic $P$,
\[
  \psi\Big(\prod_{\beta\in S}\beta\Big)=\prod(\text{roots of }P)=(-1)^{M}P.\!\operatorname{coeff}_0 .
\]
Since $\psi$ is injective and fixes $\C(t)$, the identity descends to $L$: $\prod_{\beta\in
S}\beta=(-1)^{M}P.\!\operatorname{coeff}_0$ once the latter is shown to be a base--field element. This is
the content of the next, and final, computation.

\subsection{The scalar identity $\Pi=c\,t$}
Taking the $X^0$--coefficient of \eqref{eq:weier} gives, with $h(0,t):=[X^0]h\in A^{\times}$ and
$r_0=R(0)$,
\begin{equation}\label{eq:pih}
  P.\!\operatorname{coeff}_0\cdot h(0,t)\;=\;[X^0]\Phi\;=\;-t\,r_0,
\end{equation}
so $\Pi=(-1)^M P.\!\operatorname{coeff}_0$ equals $c\,t$ with $c=(-1)^{M+1}r_0$ \emph{if and only if}
$h(0,t)=1$. The unit $h(0,t)\in\C[\![t]\!]$ is normalized by $h(0,0)=1$ (from $P\equiv X^M\equiv\Phi$ mod
$t$ and cancellation of $X^M$ in the domain $\C[\![X]\!]$). It remains to show that its $t$--derivative
vanishes:
\begin{equation}\label{eq:hderiv}
  \partial_t\,h(0,t)\;=\;0\qquad\text{when } D_m=0\ \text{for all } m\ge1 .
\end{equation}
Over a field of characteristic $0$, \eqref{eq:hderiv} together with $h(0,0)=1$ forces $h(0,t)=1$, hence
$\Pi=c\,t$, hence Proposition~\ref{prop:crux}.

Identity \eqref{eq:hderiv} is the sole analytic ingredient of the entire paper, and it is a
one--line logarithmic--derivative computation once the objects are placed in a \emph{single}
completion. Work in
\[
  \mathcal A\;=\;\big(\C(\!(x)\!)\big)[\![t]\!],
\]
the $t$--power series over the Laurent field $\C(\!(x)\!)$. Here $\Phi=x^{M}-t\,R(x)$, $P$ and $h$ all
map in and are \emph{units} (their $t$--constant terms are $x^{M},x^{M},1$, all nonzero in the field
$\C(\!(x)\!)$), so $R/\Phi,\ \partial_tP/P,\ \partial_th/h$ are honest elements of $\mathcal A$.
Differentiating $\Phi=P\cdot h$ in $t$ and dividing by $\Phi$,
\[
  -\,\frac{R}{\Phi}\;=\;\frac{\partial_tP}{P}+\frac{\partial_th}{h}.
\]
Apply the $\C[\![t]\!]$--linear functional $[x^{0}]$ (the $x^0$--Laurent coefficient, a ring
homomorphism on the ``disk'' subring of series with nonnegative $x$--support). Three evaluations finish
the argument:
\begin{enumerate}[label=(\roman*)]
\item $[x^{0}](\partial_tP/P)=0$: $P$ is monic of $x$--degree $M$ and $\partial_tP$ has $x$--degree $<M$,
so $\partial_tP/P$ has strictly negative $x$--degree and no $x^0$ term (a reversed--polynomial degree
count);
\item $[x^{0}](-R/\Phi)=-\sum_{m\ge1}D_m\,t^{m-1}$: expanding $x^{M}/\Phi=(1-tR/x^{M})^{-1}$ as a
geometric series identifies the $x^0$ coefficient with the moment generating function
$\sum_m D_m t^{m}$;
\item $[x^{0}](\partial_th/h)=\partial_t h(0,t)/h(0,t)$: $h$ and $\partial_t h$ have nonnegative
$x$--support, so on that subring $[x^0]$ is the ring homomorphism $h\mapsto h(0,t)$, which commutes with
$\partial_t$.
\end{enumerate}
Combining, $-\sum_{m\ge1}D_m t^{m-1}=0+\partial_t h(0,t)/h(0,t)$. Under $D_m\equiv0$ the left side is $0$,
so $\partial_t h(0,t)=0$, proving \eqref{eq:hderiv} and completing the proof of
Proposition~\ref{prop:crux}, hence of Theorems~\ref{thm:dvdk}, \ref{thm:nc2}, and~\ref{thm:gmc2}. \qed

\begin{remark}
The passage between the ``space'' completion $(\C(\!(x)\!))[\![t]\!]$ used for \eqref{eq:hderiv} and the
``deformation'' completion $\C[\![t]\!][X]$ used for Lemma~\ref{lem:dvd} is a transposition of the two
power--series variables; both are algebraic and neither requires extending a valuation to the algebraic
closure of $\C(\!(t)\!)$. This is the decisive simplification over the Puiseux/Newton--polygon route: the
orbit--product contradiction (Proposition~\ref{prop:no-ct}) accepts \emph{any} packet whose product is
$c\,t$, so the packet may be defined algebraically as the roots of the Weierstrass factor $P$ rather
than by valuation.
\end{remark}

\section{The Lean formalization}\label{sec:lean}

The proof is formalized in Lean~4 against \texttt{mathlib} \cite{mathlib}. We summarize the module
structure; the correspondence to the sections above is exact.

\begin{itemize}[leftmargin=1.4em]
\item \textbf{The functional and the charge reduction.} \texttt{E}, \texttt{wt}, \texttt{charge},
\texttt{NC2}, and the reduction \texttt{gmc2\_of\_nc2} (Proposition~\ref{prop:nc2-to-gmc2}) are in
\texttt{GMC2Reduction}. The threshold $N=\deg_W Q+1$ is explicit.
\item \textbf{$\text{NC}(2)\Leftarrow\text{DvdK}(1)$.} The lowest--balanced--face/Frobenius argument of
Section~\ref{sec:nc2} is \texttt{GMC2NC2} (with the normalized specialization contradiction
\texttt{normalized\_specialization\_contradiction} and the discharged height witness), yielding
\texttt{nc2\_of\_dvdK1} and \texttt{gmc2\_of\_dvdK1}.
\item \textbf{The constant--term interface.} \texttt{DvdK1} (Theorem~\ref{thm:dvdk}),
\texttt{ChargesStraddleZero}, and the universal $\mathcal C_m(q)$ \eqref{eq:ctr}
(\texttt{constantTermRelation}) are in \texttt{GMC2DvdKInterface},
\texttt{GMC2LowestFaceExistence}, \texttt{GMC2ConstantTermRelations}; the zero--charge case and the
both--signs reduction are in \texttt{GMC2DvdKZeroCharge}.
\item \textbf{Univariate reduction.} \texttt{SinglePolyCrux} (Proposition~\ref{prop:crux}), the shift
$R=\texttt{shiftedPolynomial}$, the identity \eqref{eq:checkA} (\texttt{GMC2LaurentShiftCheckA}), and the
capstone \texttt{gmc2\_of\_crux} are in \texttt{GMC2DvdKUnivariateReduction}.
\item \textbf{Orbit--product.} Lemma~\ref{lem:orbit} and Proposition~\ref{prop:no-ct} are
\texttt{GMC2OrbitProduct} and \texttt{GMC2Thm2067Wrapper}
(\texttt{thm2067\_contradiction}); irreducibility (Lemma~\ref{lem:irred}) is
\texttt{GMC2PhiIrreducible}/\texttt{GMC2DvdKAssembly}; Vieta ($\prod\alpha=d$) is
\texttt{GMC2PhiVieta} (\texttt{prod\_rootSet\_Phi}); the char--$0$ specialization discharging
separability and Galois--fixedness is \texttt{GMC2Thm2067HSonly} (\texttt{thm2067\_reduced\_to\_hS}).
\item \textbf{The frame bridge.} The valuation--free packet transport of
Section~\ref{sec:analytic} --- root transport, the packet bijection, and the assembly
$\prod_{\beta\in S}\beta=\text{algebraMap}(c\,t)$ --- is
\texttt{GMC2FrameBridgeRoots}, \texttt{GMC2FrameBridgePacket} (\texttt{exists\_packet\_prod\_eq}, a
\texttt{Finset.prod\_bij}), \texttt{GMC2FrameBridge}, \texttt{GMC2FrameBridgeAssembly}.
\item \textbf{Weierstrass and the analytic core.} \eqref{eq:weier}, $P=\texttt{smallRootFactor}$, and
\eqref{eq:pih} are \texttt{GMC2DvdKWeierstrass} (built on \texttt{mathlib}'s Weierstrass preparation);
Lemma~\ref{lem:dvd} is \texttt{GMC2FrameBridgeDvd} (\texttt{smallRootFactor\_dvd\_PhiPoly});
the normalization $h(0,0)=1$ is \texttt{GMC2DvdKUnitOrigin}; the log--derivative identity
\eqref{eq:hderiv} is proved in the unified frame $(\C(\!(x)\!))[\![t]\!]$ across
\texttt{GMC2DvdKFrame}, \texttt{GMC2DvdKTranspose}, \texttt{GMC2DvdKHderivAssembly},
\texttt{GMC2DvdKTransposeAssembly}; the char--$0$ ``vanishing derivative $\Rightarrow$ constant'' step is
\texttt{GMC2DvdKCharZeroClosing}.
\item \textbf{Discharge of the crux and the front door.} The wiring over
$\Omega=\overline{\C(\!(x)\!)}$ that proves \texttt{singlePolyCrux\_holds}, and the unconditional theorem
$\texttt{gmc2\_unconditional}=\texttt{gmc2\_of\_crux}\ \texttt{singlePolyCrux\_holds}$, are in
\texttt{GMC2DvdKOmegaWiring}; the documented statement \texttt{GMC2.gmc2} is in \texttt{GMC2Main}.
\end{itemize}

\subsection*{Verification}
The build target \texttt{GMC2Main} compiles the entire proof and is self--contained. The kernel
axiom audit
\[
  \texttt{\#print axioms GMC2.gmc2}\;=\;[\,\mathtt{propext},\ \mathtt{Classical.choice},\ \mathtt{Quot.sound}\,]
\]
certifies that the theorem uses only \texttt{mathlib}'s three foundational axioms: there is no
\texttt{sorry}, no \texttt{native\_decide} (hence no unaudited kernel reduction), and no additional
axiom. Every module terminates in a \texttt{\#print axioms} directive with the same result, so the
property holds throughout the dependency graph, not merely at the root.

\section{Concluding remarks}\label{sec:remarks}

\subsection*{Sharpness}
Theorem~\ref{thm:gmc2} is dimension--optimal: by \cite{DerksenEssenZhao2017} the Gaussian--Moments
phenomenon fails once enough Gaussian variables are present. The obstruction to higher dimensions is
genuinely multi--radial: the single hyperbolic charge $\ch=s_0-s_1$ of $\C[Z,W]$ is replaced by a
lattice of charges, and one--sidedness is no longer forced. Theorem~\ref{thm:nc2} is the exact statement
that survives in dimension two: instability of the moment functional (all powers null) coincides with
algebraic instability (a strict one--sided charge, i.e. a destabilizing one--parameter subgroup of the
hyperbolic torus).

\subsection*{The one--variable input}
The argument isolates the entire transcendental difficulty of the problem in the single scalar identity
\eqref{eq:hderiv}, a formal logarithmic derivative. This is a constructive counterpart to the analytic
proof of Duistermaat and van der Kallen \cite{DvdK1998}, replacing residues and monodromy of algebraic
functions by a Weierstrass factorization and a Galois orbit--product, and in particular giving a proof
of Theorem~\ref{thm:dvdk} that is elementary enough to be checked by a proof assistant.

\subsection*{Relation to the Jacobian Conjecture}
By \cite{DerksenEssenZhao2017}, the Gaussian--Moments Conjecture implies the Jacobian Conjecture. Since
that conjecture fails in higher dimension, the implication cannot be run in full; the two--dimensional
case established here is its sharp positive instance.

\subsection*{Availability}
The Lean sources build against \texttt{mathlib} and reproduce the axiom audit above. The proof was
developed and machine--checked in the \texttt{TournamentH7} Lean project; the theorem \texttt{GMC2.gmc2}
and its self--contained build target \texttt{GMC2Main} are the canonical entry points.

\begin{thebibliography}{9}

\bibitem{DerksenEssenZhao2017}
H.~Derksen, A.~van den Essen, and W.~Zhao,
\emph{The Gaussian moments conjecture and the Jacobian conjecture},
Israel J. Math. \textbf{219} (2017), no.~2, 917--928. \texttt{arXiv:1506.05192}.

\bibitem{DvdK1998}
J.~J.~Duistermaat and W.~van der Kallen,
\emph{Constant terms in powers of a Laurent polynomial},
Indag. Math. (N.S.) \textbf{9} (1998), no.~2, 221--231.

\bibitem{Mathieu1997}
O.~Mathieu,
\emph{Some conjectures about invariant theory and their applications},
in \emph{Alg\`ebre non commutative, groupes quantiques et invariants}
(Reims, 1995), S\'emin. Congr. \textbf{2}, Soc. Math. France, 1997, pp.~263--279.

\bibitem{Zhao2010}
W.~Zhao,
\emph{Generalizations of the image conjecture and the Mathieu conjecture},
J. Pure Appl. Algebra \textbf{214} (2010), no.~7, 1200--1216.

\bibitem{Zhao2007a}
W.~Zhao,
\emph{Hessian nilpotent polynomials and the Jacobian conjecture},
Trans. Amer. Math. Soc. \textbf{359} (2007), no.~1, 249--274.

\bibitem{ZhaoImage}
W.~Zhao,
\emph{Images of commuting differential operators of order one with constant leading coefficients},
J. Algebra \textbf{324} (2010), no.~2, 231--247.

\bibitem{vandenEssen2000}
A.~van den Essen,
\emph{Polynomial Automorphisms and the Jacobian Conjecture},
Progress in Mathematics \textbf{190}, Birkh\"auser, Basel, 2000.

\bibitem{mathlib}
The mathlib Community,
\emph{The Lean mathematical library},
in \emph{CPP 2020}, ACM, 2020, pp.~367--381.

\end{thebibliography}

\end{document}
